\subsection{A relative sign obstruction for more source monoids}

\begin{proposition}[Proved: row-independent actions over permutation-row sources]
\label{CP2:A-row-independent-permutation-source}
Let $M$ have minimal ideal
$K=\{x_{ij}:i,j\in\{0,1\}\}$, a rectangular band.
Suppose every element outside $K$ acts on its two minimal
rows by a permutation.  Write $T=M\setminus K$ and let
$\chi:T\to\{1,-1\}$ be the sign of that row permutation.
Assume the one-dimensional left $\mathbb QT$-module
$\mathbb Q_\chi$ is projective.

Let $S$ be a left-$N$, right-$M$ biset with no globally
right-fixed state and with
\[
 sx_{0j}=sx_{1j}\qquad(s\in S,\ j=0,1).
\]
If $|N|\le7$, then every rational degree-two square of
the collage is zero whenever its second rational
cohomology has dimension at most one.
More generally, the same conclusion holds if the ordinary
left $N$-act cone on the set of right $K$-core orbits has
zero degree-two products.
\end{proposition}
\begin{proof}
The outside set $T$ is a submonoid: a product of two
outside actors still permutes the rows, whereas every
band actor is constant on them.  Thus
$A=\mathbb QM/\mathbb QK$ is exactly $\mathbb QT$.
Interchange the two rows of $K$ and fix every element of
$T$.  This defines a monoid automorphism $\sigma$.
Indeed, row interchange commutes with each outside row
permutation; right multiplication only changes columns;
and products in $T$ stay in $T$.

The core $U=SK$ is a disjoint union of regular two-state
column orbits.  A one-state core would be fixed by the
whole source, since $K$ is an ideal.  Every element of
$M$ preserves each core orbit and acts there by its
column transformation.  Every commuting $N$ map preserves
the column labels and acts on the orbit set $B$.
Consequently $U=B\times eM$ as a biset, for
$e=x_{00}$, with trivial $M$ action on $B$.
The regular-core action bar is contractible, and its
projection to $BM$ has an $N$-equivariant natural
nullhomotopy: on state $(b,j)$ use the arrow $x_{0j}$;
the identity $x_{0j}m=x_{0,\tau_m(j)}$ verifies naturality.
Therefore the actual core subcollage has the splitting
\[
 X_U\simeq BM\vee
 \operatorname{hocofib}\bigl(B(N,B)\longrightarrow BN\bigr).
\]
Every degree-two product on $BM$ is zero by the
two-row conclusion of Corollary~\ref{A:monoids-nine}
and rational polarization.  The displayed hypothesis on the ordinary
act cone makes all such products on $X_U$ zero as well.

Put $R=\mathbb QM$, and let $V$ be the row augmentation
module in the exact sequence
\[
 0\longrightarrow V\longrightarrow Re
   \longrightarrow\mathbb Q\longrightarrow0.
\]
It is the inflated $A$-module $\mathbb Q_\chi$.
The semilinear row interchange acts as $-1$ on $V$.
The quotient module $L=\mathbb Q[S]/\mathbb Q[U]$ is
annihilated by $K$, and hence $Le=0$.  Tensoring this
exact sequence derivedly gives a natural comparison
\[
 L\otimes_R^{\mathbf L}\mathbb Q
  \simeq (L\otimes_R^{\mathbf L}V)[1]
  \simeq (L\otimes_A^{\mathbf L}V)[1]
  \simeq (L\otimes_A V)[1].
\]
For the middle comparison, the ideal factorization
$\mathbb QK=Re\otimes_{\mathbb Q}eR$ makes
$0\to\mathbb QK\to R\to A\to0$ a left-projective
resolution, and $Le=0$ kills its positive tensor term.
The last comparison uses precisely the projectivity of
$V$ over $A$.

Row independence makes the identity on $S$, together
with $\sigma$ on $M$, an involution of the entire
collage and its core subcollage.  It is the identity on
$L$ and on $A$, and acts as $-1$ on $V$.
The preceding natural relative comparison therefore has
sign $-1$.  It is $N$-equivariant, so taking derived
$N$ coinvariants preserves both the comparison and its
sign.  Finally, the actual mapping-cylinder pair gives
$X/X_U\simeq\Sigma(E/E_U)$, adding only a degree shift.
Thus the involution acts by $-1$ on every rational
relative homology group $H_*(X,X_U)$.

Exactness of rational involution eigenspaces now makes
$H_4(X_U)\to H_4(X)^+$ surjective.  When $H^2(X)$ is
one-dimensional, the square of its generator is invariant
under the involution.  It pairs to zero with the negative
fourth-homology eigenspace by its sign, and with the
positive eigenspace because that eigenspace comes from
$X_U$, where products vanish.  The square is zero.
The zero-dimensional second-cohomology case is immediate.
For $|N|\le7$, the required ordinary-act product
vanishing is Corollary~\ref{CP2:A-high-seven-act}.
\end{proof}

\begin{lemma}[Proved: a literal projectivity certificate]
\label{CP2:A-permutation-source-projective-character}
In the preceding proposition, it suffices to exhibit an
idempotent $f\in T$ such that $G=fTf$ is a group with
identity $f$ and $Tf=G$.
\end{lemma}
\begin{proof}
Since $\chi(f)=1$, the element
\[
 p=\frac1{|G|}\sum_{g\in G}\chi(g)g\in\mathbb QT
\]
satisfies $tp=\chi(t)p$ for all $t\in T$.
Indeed $tf\in G$, multiplication by $tf$ permutes $G$,
and $\chi(tf)=\chi(t)$.  Also $\chi(p)=1$, so $p^2=p$.
Thus $\mathbb QT p\cong\mathbb Q_\chi$ is a direct
summand of the regular module.  These group, closure and
signed eigenvector identities are finite literal table
equations and supply a separate checkable witness.
\end{proof}

\begin{remark*}[Scope before using the new gate]
This proposition supplies an additional analytic omission.
It does not change any earlier frozen enumeration or
promote a partial stream to a completed scope.  A search
using it must identify every omitted source/right-action
pair, verify the row-independence and projectivity witness,
and account separately for all previously emitted records.
\end{remark*}
